The expression being summed on the left-hand side differs from the sum in (108.2) by the omission of one factor (a $3j$-symbol). Thus the sum in (108.4) may be represented by a tetrahedron (Fig.~45) with one face removed; this is why the sum is not a scalar. In other words, in its transfor\SourcePageJoin{536}{539}mation properties it corresponds to a single $3j$-symbol---the one on the right-hand side of (108.4)---and must therefore be proportional to that symbol. The proportionality factor (the $6j$-symbol on the right-hand side) is obtained at once by multiplying both sides by
\[
\begin{pmatrix}j_1&j_2&j_3\\m_1&m_2&m_3\end{pmatrix}
\]
and summing over the remaining numbers $m_1,m_2,m_3$.
